\section{Advanced Variants of Grover's Algorithm}
\label{sec:advanced_grover}
%
%

As mentioned in the previous section, the standard Grover's algorithm is  optimal for general unstructured search, but it has a problem in it. 
It's success probability is not exactly 100\%, as also, what will be a bigger issue, it has the overshooting effect. 
To address these issues several other variants of the Grover's algorithm have been developed.
%
%

\subsection{Exact Quantum Search}
%
%

\subsubsection{Idea of Exact Quantum Search}

In the standard version of Grover's algorithm we perform $k \approx \frac{\pi}{4}\sqrt{\frac{N}{M}}$ iterations, where $M$ is the number of solutions. 
Let's further denote $\sqrt{\frac{N}{M}}$ as $Q$
As we know $k$ must be an integer which mean it already an apporximation so we may not align perfectly with the correct angle. 
Also due to the formula using small angle approximation, the number of iterations may differ even more from the target state vector $\ket{x^*}$. 
The final success probability $P_k = \sin^2((2k+1)\theta)$ is typically 
quite close to $1$, but sometimes that may result in an issue.
%
%

\bigskip
The \textit{Exact Quantum Search} variant, first proposed by Long \cite{long2001grover}, modifies the algorithm 
to guarantee correct answer (excluding the errors from floating point arithmetic). 
In short, it peforms the standard Grover iterations until the state vector is within a single step of the target, 
then using one step it will try to align with the target angle perfectly.
%
%

\bigskip
In the standard algorithm, every iteration rotates the state vector by an angle of $2\theta$. 
After $k$ full iterations, the total angle relative to the starting uniform superposition is $(2k+1)\theta$.
As described earlier, to perform a correct measurement, that angle should equal to $\frac{\pi}{2}$. 
So to solve the problem with $k$ being a whole number, we should at the end rotate the vector by $\delta$, which is defined as $\delta = \frac{\pi}{2} - (2k+1)\theta$.
%
%

\bigskip
To also ensure that this computation is fully accurate, we must stop relying on the small angle approximation (which was $\theta \approx \frac{1}{Q}$). 
Instead we will calculate the initial angle using a formula:
\bigskip
\begin{equation}
    \theta = \arcsin (\frac{1}{Q})
\end{equation}
\bigskip
Due to that the gap $\delta$ is known precisely.
%
%

\bigskip
The inital angle formula comes from the definition of the uniform superposition. 
At the start of the algorithm, the probability amplitude of finding the target state is exactly $\frac{1}{Q}$ 
(comes right from the distribution after using the Hadamard gate). 
%
%
In the geometric representation, this initial amplitude corresponds to the sine of the starting angle $\theta$, 
which gives us the exact equation $\sin(\theta) = \frac{1}{Q}$. 
By applying the inverse sine function, we can extract the precise angle.
%
%

\bigskip
Using that two improvements we can in theory achieve the deterministic result of the search algorithm.

\subsubsection{Implementation of Exact Quantum Search}

To implement this variant we must implement custom parameterized diffusion operator to apply specific fractional phase shifts, which in that case will be the exact 
remaining angle.

\bigskip
The code below shows the implementation of the Exact Quantum Search. 
The algorithm performs standard iterations first. Then it calculates the remaining angle and applies generated diffuser for the final fractional step.
%
%

%After changes it may not be needed
\newpage

\begin{lstlisting}[
    language=Python, 
    title={Implementation of the Exact Quantum Search algorithm.}, 
    label={lst:exact_impl},              
    basicstyle=\ttfamily\small,    
    backgroundcolor=\color{gray!5},
    frame=single, 
    rulecolor=\color{gray!40},     
    framesep=5pt,                 
    numbers=left, 
    numberstyle=\tiny\color{gray}, 
    numbersep=10pt,                
    commentstyle=\color{green!50!black}\itshape, 
    keywordstyle=\color{blue}\bfseries, 
    stringstyle=\color{red!80!black},
    showstringspaces=false,
    breaklines=true,
    aboveskip=1.5em,              
    belowskip=1.5em                
]
class GroverAlgorithmExact(GroverAlgorithmBase):
    def __init__(self, oracle):
        super().__init__(oracle)
        self.standard_diffuser = self._build_diffuser(np.pi)

    def _build_diffuser(self, angle: float):
        qc = QuantumCircuit(self.num_qubits, name=f"U_s")
        qc.h(range(self.num_qubits))
        qc.x(range(self.num_qubits))

        if self.num_qubits > 1:
            mcphase = MCPhaseGate(angle, self.num_qubits - 1)
            qc.append(mcphase, range(self.num_qubits))
        else:
            qc.p(angle, 0)

        qc.x(range(self.num_qubits))
        qc.h(range(self.num_qubits))
        return qc

    def build_circuit(self, num_solutions: int = None):
        if number_of_solutions is None or number_of_solutions < 1:
            raise ValueError("Invalid number of solutions.")

        n_states = 2 ** self.num_qubits
        theta = 2 * np.arcsin(np.sqrt(num_solutions / n_states))
        exact_iters_float = (np.pi / (2 * theta)) - 0.5
        iterations = max(0, math.floor(exact_iters_float))

        qc = QuantumCircuit(self.num_qubits, self.num_qubits)
        qc.h(range(self.num_qubits))

        for _ in range(iterations):
            qc.compose(self.oracle, inplace=True)
            qc.compose(self.standard_diffuser, inplace=True)

        # Perform the final fractional iteration
        remaining_angle = np.pi/2 - (iterations + 0.5) * theta
        if remaining_angle > 1e-5:
            final_phase = 2 * remaining_angle 
            final_diffuser = self._build_diffuser(final_phase)
            qc.compose(self.oracle, inplace=True)
            qc.compose(final_diffuser, inplace=True)

        qc.measure(range(self.num_qubits), range(self.num_qubits))
        return qc
\end{lstlisting}
%
%

\subsection{Fixed-Point Quantum Search}
%
%
\subsubsection{Idea of Fixed-Point Quantum Search}

A significant problem of the standard algorithm is its periodic. 
If the number of target solutions $M$ is unknown, it is impossible (or at least hard) to compute the optimal number of iterations $k$ to perform. 
Performing too may iterations in the algorithm causes the state vector to rotate past the $\ket{x^*}$ axis, which results in
decreasing the success probability (Overshooting effect).
%
%

\bigskip
The \textit{Fixed-Point Quantum Search} addresses this by changing the geometry of the operations. 
Introduced by Yoder, Low, and Chuang \cite{yoder2014fixed}, this variant replaces the standard $\pi$ ($180^\circ$) 
phase shifts in both the phase oracle ($U_f$) and the diffusion operator ($U_s$) with a phase shift of $\pi/3$ ($60^\circ$).
%
%


\subsubsection{Proof of Convergence}
%
%

To prove that replacing the standard phase shifts to $\pi/3$ will not be periodic but convergent, we must analyze failure probability amplitude and prove that it is 
strictly decreasing. 
Let the initial uniform superposition state be defined as $\ket{s} = \sin\gamma\ket{x^*} + \cos\gamma\ket{s'}$, 
where $\ket{s'}$ is the uniform superposition of all states different than a target one and $\gamma$ is the initial angle.
%
%

\bigskip
As said before in this variant we will use phase shift of $\omega = e^{i\pi/3}$. It follows that the modified oracle $U_f(\pi/3)$ and diffuser $U_s(\pi/3)$ operators 
can be denoted as:
\bigskip
\begin{equation}
    U_f = I - (1 - \omega)\ket{x^*}\bra{x^*}
\end{equation}
\bigskip
\begin{equation}
    U_s = I - (1 - \omega)\ket{s}\bra{s}
\end{equation}
%
%
\bigskip

Applying the oracle to the initial state yields $U_f \ket{s} = \omega\sin\gamma\ket{x^*} + \cos\gamma\ket{s'}$. 
Also applying a one full new iteration gives us the following state:
\bigskip
\begin{equation}
    \mathcal{G}_{\pi/3} \ket{s} = U_s U_f \ket{s} = (I - (1 - \omega)\ket{s}\bra{s}) (U_f \ket{s}) = U_f \ket{s} - (1 - \omega)\ket{s} \bra{s} U_f \ket{s}
\end{equation}
%
%
\bigskip

Let's observe the amplitude of the state different than the target one. 
This amplitude will represent the error of our algorithm and can be denoted as $\bra{s'} U_s U_f \ket{s}$:
\bigskip
\begin{equation}
    \bra{s'} U_s U_f \ket{s} = \bra{s'} U_f \ket{s} - (1 - \omega)\bra{s'} \ket{s} \bra{s} U_f \ket{s}
\end{equation}
%
%
\bigskip

Using a few substitutions form the definition of $\ket{s}$ and $U_f \ket{s}$ we can evaluate a few terms: 
\bigskip
\begin{equation}
    \bra{s'} \ket{s} = \sin \gamma \bra{s'}\ket{x^*} + \cos \gamma \bra{s'} \ket{s'} = \sin \gamma \cdot 0 +  \cos \gamma \cdot 1 = \cos \gamma
\end{equation}
\bigskip
\begin{equation}
    \bra{s'} U_f \ket{s} = \omega \sin \gamma \bra{s'} \ket{x^*} + \cos \gamma \bra{s'} \ket{s'} = \omega \sin \gamma \cdot 0 +  \cos \gamma \cdot 1 = \cos \gamma
\end{equation}
\bigskip
\begin{equation}
    \bra{s} U_f \ket{s} = (\sin\gamma\bra{x^*} + \cos\gamma\bra{s'})(\omega\sin\gamma\ket{x^*} + \cos\gamma\ket{s'}) = \omega\sin^2\gamma + \cos^2\gamma
\end{equation}
\bigskip

Using all that we obtain that our total error can be represented as:
\bigskip
\begin{equation}
    \bra{s'} U_s U_f \ket{s} = \cos\gamma - (1 - \omega)\cos\gamma (\omega\sin^2\gamma + \cos^2\gamma)
\end{equation}
%
%
\bigskip

Factoring out $\cos\gamma$ simplifies the expression to:
\bigskip
\begin{equation}
    \cos\gamma [ 1 - (\omega - \omega^2)\sin^2\gamma - (1 - \omega)\cos^2\gamma ]
\end{equation}
%
%
\bigskip

We know that $\omega = e^{i\pi/3}$ is a root of the equation $\omega^2 - \omega + 1 = 0$, 
which implies that $\omega - \omega^2 = 1$. Substituting this into the above equation gives us:
\begin{equation}
    \cos\gamma [ 1 - \sin^2\gamma - (1 - \omega)\cos^2\gamma ] = \cos\gamma [ \cos^2\gamma - (1 - \omega)\cos^2\gamma ] = \omega \cos^3\gamma
\end{equation}
\bigskip
%
%

The value of amplitude for non target elements after one iteration equals to $\cos^3\gamma$. 
As we know that the initial angle $\gamma > 0$ (as the at least one target element exists), we know that $\cos\gamma < 1$.
%
%

\bigskip
After $k$ iterations, the error amplitude can be bounded by $\cos^{2k+1}\gamma$, which now can be clearly seen approaches $0$ 
as $k \to \infty$. That means that the probability of measuring the correct state approaches $1$ and the overshooting effect would not have place there.
%
%

\subsubsection{Implementation of Fixed-Point Quantum Search}

The implementation is quite similar to the oryginal Grover implementation. 
The main difference here is to create oracle and diffuser $U_s(\pi/3)$, which will be applyed as in the classic Grover's algorithm. The problem here in the general 
will be that $U_f(\pi/3)$ cannot be created without a knowledge about the black box and has to be given as an input. In practice Qiskit can handle that as it knows how that 
oracle is constructed but it may not always be such easy case.
%
%

% May not be needed later
\newpage

\begin{lstlisting}[
    language=Python, 
    title={Implementation of the Fixed-Point Quantum Search algorithm.}, 
    label={lst:fixed_impl},              
    basicstyle=\ttfamily\small,    
    backgroundcolor=\color{gray!5},
    frame=single, 
    rulecolor=\color{gray!40},     
    framesep=5pt,                 
    numbers=left, 
    numberstyle=\tiny\color{gray}, 
    numbersep=10pt,                
    commentstyle=\color{green!50!black}\itshape, 
    keywordstyle=\color{blue}\bfseries, 
    stringstyle=\color{red!80!black},
    showstringspaces=false,
    breaklines=true,
    aboveskip=1.5em,              
    belowskip=1.5em                
]
class GroverAlgorithmFixedPoint(GroverAlgorithmBase):
    def __init__(self, oracle):
        super().__init__(oracle)

        self.diffuser_pi_3 = self._build_diffuser(np.pi / 3)
        self.oracle_pi_3 = self.oracle.power(1/3)

    def _build_diffuser(self, angle: float) -> QuantumCircuit:
        qc = QuantumCircuit(self.num_qubits, name=f"U_s")
        qc.h(range(self.num_qubits))
        qc.x(range(self.num_qubits))

        if self.num_qubits > 1:
            mcphase = MCPhaseGate(angle, self.num_qubits - 1)
            qc.append(mcphase, range(self.num_qubits))
        else:
            qc.p(angle, 0)

        qc.x(range(self.num_qubits))
        qc.h(range(self.num_qubits))
        return qc

    def build_circuit(self,  num_solutions: int = None):
        opt_its = self._calculate_optimal_iterations(num_solutions)

        # It will converge securely without overshooting
        iterations = 2*(opt_its) + 17

        qc = QuantumCircuit(self.num_qubits, self.num_qubits)
        qc.h(range(self.num_qubits))

        # Iterations converge securely without overshooting
        for _ in range(iterations):
            qc.compose(self.oracle_pi_3, inplace=True)
            qc.compose(self.diffuser_pi_3, inplace=True)

        qc.measure(range(self.num_qubits), range(self.num_qubits))
        return qc
\end{lstlisting}